
\clearpage
\section*{Deleted original main-article tables}
The following twelve tables reproduce the removed original tables in full for
comparison with the replacement tables above. Red text marks material deleted
from the original submission. These numerical values and methodological accounts
are historical and do not describe the revised analysis. Table numbers below
refer to the original submission. This deletion record appears only in the
marked manuscript; it is not part of the clean revised article.
\begingroup
\singlespacing
\color{red!70!black}
\setcounter{table}{0}
\renewcommand{\thetable}{Original \arabic{table}}
\renewcommand{\theHtable}{deletedoriginal.\arabic{table}}

\begin{table}[H]\color{red!70!black}
\centering
\caption{Null-model validation statistics for whole-brain coevolutionary energy. After unified
outlier removal, 78 ASD and 86 TD participants remained. For each subject, empirical energy was
compared to the mean energy from topology-preserving null networks in which node states and edge
signs were randomly permuted. Shapiro--Wilk tests confirmed normality of the difference
distributions, justifying the use of paired $t$-tests alongside Wilcoxon signed-rank tests.}
\label{deletedoriginal:tab:null_model_validation}
\small
\resizebox{\linewidth}{!}{%
\begin{tabular}{lccccccc}
\toprule
Group & $N$ & Shapiro--Wilk $W$ & Shapiro--Wilk $p$ & $t$ (df) & $t$-test $p$ & Wilcoxon $W$ & Wilcoxon $p$ \\
\midrule
ASD & 90 & 0.992 & 0.87 & $-4.12$ (89) & $7.5\times 10^{-5}$ & 1024 & $9.1\times 10^{-5}$ \\
TD  & 92 & 0.989 & 0.81 & $-3.98$ (91) & $1.2\times 10^{-4}$ & 1008 & $1.6\times 10^{-4}$ \\
\bottomrule
\end{tabular}
}
\end{table}
\clearpage

\begin{table}[H]\color{red!70!black}
\centering
\caption{Summary of statistical procedures and group comparison tests for whole-brain coevolutionary
metrics. The upper panel lists the main analysis steps and their purpose. The lower panel reports
group comparison statistics (Mann--Whitney $U$ tests and independent $t$-tests where appropriate,
selected based on Shapiro--Wilk normality tests and Levene's homogeneity test applied
independently to each feature), raw $p$-values, and FDR-corrected $p$-values for ASD versus TD
after ComBat harmonization (GSR pipeline, $n = 78$ ASD, $86$ TD). No whole-brain feature survived
FDR correction.}
\label{deletedoriginal:tab:assumptions_and_tests}
\small
\setlength{\tabcolsep}{3pt}
\resizebox{\linewidth}{!}{%
\begin{tabular}{p{0.32\columnwidth}p{0.30\columnwidth}p{0.23\columnwidth}p{0.13\columnwidth}}
\toprule
Procedure & Purpose & Method / Test & Threshold / Correction \\
\midrule
ComBat harmonization  & Remove multi-site batch effects while preserving biological signal & Per-feature empirical Bayes adjustment & --- \\
Null-model generation     & Destroy fALFF--FC alignment while preserving graph topology & 1\,000 permutations of node states and edge signs & --- \\
Outlier removal           & Exclude extreme values in null distributions and empirical metrics & Tukey's IQR criterion ($\pm 1.5\times\mathrm{IQR}$) & --- \\
Normality assessment      & Guide parametric vs. non-parametric test selection & Shapiro--Wilk test & $p > 0.05$ (normal) \\
Variance assessment       & Distinguish standard $t$-test from Welch's & Levene's test & $p > 0.05$ (equal variances) \\
Parametric comparison     & Test mean differences between ASD and TD & Independent-samples $t$-test or Welch's & two-tailed, $\alpha = 0.05$ \\
Nonparametric comparison  & Applied when normality is violated in either group & Mann--Whitney $U$ test & two-tailed, $\alpha = 0.05$ \\
Multiple comparisons      & Control false discovery across families of tests & Benjamini--Hochberg FDR & $q < 0.05$ \\
\bottomrule
\end{tabular}
}

\vspace{0.8em}

\resizebox{\columnwidth}{!}{%
\resizebox{\linewidth}{!}{%
\begin{tabular}{lccccc}
\toprule
Metric & ASD Mean (SD) & TD Mean (SD) & Test & $p$ & $p_{\mathrm{FDR}}$ \\
\midrule
Energy global           & $-139.7$ (78.7) & $-118.8$ (76.2)  & Mann--Whitney $U$ & $0.048$ & $0.19$ \\
Prop imbalanced same    & $0.256$ (0.006)  & $0.258$ (0.005)  & $t$-test         & $0.047$ & $0.19$ \\
Prop agreement          & $0.241$ (0.005)  & $0.240$ (0.005)  & $t$-test          & $0.070$ & $0.19$ \\
Prop disagreement       & $0.267$ (0.005)  & $0.267$ (0.005)  & $t$-test         & $0.679$ & $0.998$ \\
Prop imbalanced opp     & $0.236$ (0.005)  & $0.236$ (0.005)  & $t$-test         & $0.816$ & $0.998$ \\
Bipolarity              & $0.626$ (0.017)  & $0.627$ (0.015)  & $t$-test          & $0.532$ & $0.998$ \\
Global efficiency       & $0.283$ (0.019)  & $0.283$ (0.018)  & $t$-test          & $0.998$ & $0.998$ \\
\bottomrule
\end{tabular}
}%
}
\end{table}
\clearpage

\begin{table}[H]\color{red!70!black}
\centering
\small
\caption{Largest intra-network trends within the seven canonical Yeo resting-state networks under
the harmonized GSR pipeline ($n = 78$ ASD, $86$ TD). For each metric, the table reports ASD and TD
means, the $p$-value from the adaptive group comparison (independent-samples $t$-test or
Mann--Whitney $U$, selected by Shapiro--Wilk normality testing), and the Benjamini--Hochberg
FDR-corrected value within the intra-network family. \emph{No intra-network metric reached even
nominal significance} ($p < 0.05$) after harmonization; values are listed for completeness.}
\label{deletedoriginal:tab:network_results}
\setlength{\tabcolsep}{5pt}
\resizebox{\linewidth}{!}{%
\begin{tabular}{llcccc}
\toprule
\textbf{Network} & \textbf{Metric} & \textbf{ASD mean} & \textbf{TD mean} & \textbf{$p$} & \textbf{$p_{\mathrm{FDR}}$} \\
\midrule
Salience/VentAttn & AvgWeight   & $-0.001$ & $0.002$  & $0.064$ & $0.566$ \\
SomatoMotor       & EnergyBlock & $-1.41$  & $-0.53$  & $0.108$ & $0.566$ \\
Limbic            & EnergyBlock & $-1.13$  & $-0.29$  & $0.137$ & $0.566$ \\
SomatoMotor       & AvgFalffW   & $0.005$  & $0.002$  & $0.142$ & $0.566$ \\
DorsalAttn        & EnergyBlock & $-1.21$  & $-0.44$  & $0.163$ & $0.566$ \\
\bottomrule
\end{tabular}
}
\end{table}
\clearpage

\begin{table}[H]\color{red!70!black}
\centering
\caption{Inter-network differences in average connectivity (AvgWeight), fALFF-aligned connectivity
(AvgFalffW), and coevolutionary energy (EnergyBlock) between canonical resting-state networks. For
each network pair, group comparisons were performed on ASD versus TD using the adaptive statistical
pipeline (test selected based on Shapiro--Wilk normality assessment per feature), and
Benjamini--Hochberg FDR correction was applied across all 63 inter-network metrics. The table lists
the metrics with the smallest uncorrected $p$-values ($n = 78$ ASD, $86$ TD). No inter-network
metric survived FDR correction at $\alpha = 0.05$; the SomatoMotor--Frontoparietal effect
($p_{\mathrm{FDR}} = 0.069$) is the closest and is treated as exploratory.}
\label{deletedoriginal:tab:internetwork_new}
\resizebox{\linewidth}{!}{%
\begin{tabular}{llccc}
\toprule
\textbf{Network Pair} & \textbf{Metric} & \textbf{$p$} & \textbf{$p_{\mathrm{FDR}}$} \\
\midrule
SomatoMotor $\leftrightarrow$ Frontoparietal & AvgWeight & $0.0011$ & $0.069$ \\
SomatoMotor $\leftrightarrow$ Salience/VentAttn & AvgWeight & $0.0030$ & $0.095$ \\
DorsalAttn $\leftrightarrow$ Salience/VentAttn & AvgWeight & $0.0076$ & $0.160$ \\
DorsalAttn $\leftrightarrow$ Frontoparietal & AvgWeight & $0.015$ & $0.237$ \\
Visual $\leftrightarrow$ Default & EnergyBlock & $0.047$ & $0.590$ \\
\bottomrule
\end{tabular}
}
\end{table}
\clearpage

\begin{table}[H]\color{red!70!black}
\centering
\small
\caption{Partial Least Squares regression cross-validation results (harmonized GSR pipeline). For
each clinical score, the table reports the number of ASD subjects with valid data and the mean
cross-validated $R^2$ for the single-component model. All values are negative, indicating that the
models do not predict behavioral scores above chance---consistent with the univariate result that no
feature--clinical score association survives FDR correction at $q < 0.05$, and with the underdetermined
nature of the problem (90 features, $n = 29$--59).}
\label{deletedoriginal:tab:pls_results}
\resizebox{\linewidth}{!}{%
\begin{tabular}{lcc}
\toprule
Clinical Score & $n$ & CV $R^2$ (1 component) \\
\midrule
ADOS Total          & 59 & $-0.40$ \\
ADOS Communication  & 50 & $-0.79$ \\
ADI-R Social        & 32 & $-0.96$ \\
ADI-R Verbal        & 32 & $-0.86$ \\
ADI-R Onset         & 29 & $-0.80$ \\
\bottomrule
\end{tabular}
}
\end{table}
\clearpage

\begin{table}[H]\color{red!70!black}
\centering
\small
\caption{Strongest associations between network metrics and clinical scores under the harmonized GSR
pipeline. The table reports the strongest associations (by uncorrected Pearson $p$-value) from the
full analysis of 92 features $\times$ 5 clinical scores (460 tests), together with the whole-brain
Hamiltonian for reference. None survived FDR correction at $q < 0.05$ (all $p_{\mathrm{FDR}} > 0.59$).
The number of nominally significant associations ($\approx 21$) is close to the chance expectation.}
\label{deletedoriginal:tab:behaviour_corr_expanded}
\scriptsize
\resizebox{\linewidth}{!}{%
\begin{tabular}{p{4.6cm}lccccc}
\toprule
\textbf{Network Feature} & \textbf{Clinical Score} & \textbf{$n$} & \textbf{$r$} & \textbf{$p$} & \textbf{$p_{\mathrm{FDR}}$} & \textbf{$\rho$} \\
\midrule
Intra\_SomatoMotor\_AvgWeight & ADOS Total & 59 & $+$0.37 & 0.004 & 0.60 & $+$0.40 \\
Inter\_Vis\_\_Sal/VA\_EnergyBlock & ADI-R Social & 32 & $-$0.49 & 0.005 & 0.60 & $-$0.53 \\
Inter\_Vis\_\_Sal/VA\_AvgFalffW & ADI-R Social & 32 & $+$0.48 & 0.005 & 0.60 & $+$0.52 \\
Inter\_DA\_\_FP\_EnergyBlock & ADI-R Onset & 29 & $+$0.49 & 0.006 & 0.60 & $+$0.49 \\
Inter\_DA\_\_FP\_AvgFalffW & ADI-R Onset & 29 & $-$0.49 & 0.007 & 0.60 & $-$0.51 \\
Inter\_Vis\_\_Lim\_AvgWeight & ADOS Comm. & 50 & $-$0.34 & 0.017 & 0.90 & $-$0.33 \\
Inter\_Lim\_\_FP\_EnergyBlock & ADI-R Social & 32 & $+$0.40 & 0.022 & 0.90 & $+$0.35 \\
Whole\_Energy\_global & ADOS Total & 59 & $-$0.10 & 0.47 & 1.00 & $-$0.10 \\
\bottomrule
\end{tabular}
}
\begin{flushleft}
\footnotesize
Abbreviations: Sal/VA = Salience/VentAttn, DA = DorsalAttn, FP = Frontoparietal, Vis = Visual,
Lim = Limbic. $p_{\mathrm{FDR}}$ denotes Benjamini--Hochberg corrected Pearson $p$-values. After
harmonization, the whole-brain Hamiltonian is essentially uncorrelated with ADOS severity
($r = -0.10$), and no association survives correction.
\end{flushleft}
\end{table}
\clearpage

\begin{table}[H]\color{red!70!black}
\centering
\caption{Performance of supervised classifiers for ASD versus TD classification using the selected
coevolutionary and connectivity features after ComBat harmonization (GSR pipeline). CV accuracy was
obtained from stratified 5-fold cross-validation on the training set (80\% of participants); test
metrics were computed on the held-out 20\% test set ($n = 34$), which was not used at any stage of
feature selection or hyperparameter tuning. No model exceeded chance in cross-validation.}
\label{deletedoriginal:tab:ml_summary_full}
\setlength{\tabcolsep}{4pt}
\small
\resizebox{\linewidth}{!}{%
\begin{tabular}{lccccc}
\toprule
Model & CV Acc & Test Acc & Test Bal Acc & Test F1 & Test AUC \\
\midrule
Gaussian Naive Bayes  & 0.585 & 0.588 & 0.583 & 0.533 & 0.604 \\
SVM (RBF)             & 0.586 & 0.647 & 0.649 & 0.647 & 0.635 \\
Logistic Regression   & 0.556 & 0.500 & 0.503 & 0.514 & 0.552 \\
KNN                   & 0.594 & 0.618 & 0.608 & 0.519 & 0.639 \\
XGBoost               & 0.579 & 0.647 & 0.646 & 0.625 & 0.740 \\
\bottomrule
\end{tabular}
}
\end{table}
\clearpage

\begin{table}[H]\color{red!70!black}
\centering
\caption{Detailed metrics for the Gaussian Naive Bayes classifier (the model originally highlighted)
on the held-out test set ($n = 34$) after ComBat harmonization. Performance is near chance.
Precision, recall, and F1-score refer to the ASD class as the positive class; ROC AUC was computed
from the predicted ASD class probabilities.}
\label{deletedoriginal:tab:ml_best_model_metrics}
\setlength{\tabcolsep}{6pt}
\small
\resizebox{\linewidth}{!}{%
\begin{tabular}{lc}
\toprule
Metric & Value \\
\midrule
Accuracy & 0.588 \\
Balanced accuracy & 0.583 \\
F1-score & 0.533 \\
ROC AUC & 0.604 \\
\bottomrule
\end{tabular}
}
\end{table}
\clearpage

\begin{table}[H]\color{red!70!black}
\centering
\caption{Whole-brain coevolutionary metrics in the no-GSR sensitivity analysis (CompCor-based
preprocessing). The table reports ASD and TD group means, standard deviations, the statistical
test used (selected by Shapiro--Wilk normality criterion per feature), raw $p$-values, and
FDR-corrected $p$-values. The FDR family for the Whole-brain category includes 8 features
(6 coevolutionary metrics plus global and local graph efficiency). No whole-brain feature survived
FDR correction; bipolarity showed the smallest uncorrected $p$-value ($p = 0.026$,
$p_{\mathrm{FDR}} = 0.21$), and graph efficiency metrics
did not differ significantly between groups. Compare with Table~\ref{tab:assumptions_and_tests}
for the GSR-based primary analysis.}
\label{deletedoriginal:tab:noGSR_whole}
\small
\resizebox{\columnwidth}{!}{%
\resizebox{\linewidth}{!}{%
\begin{tabular}{lccccc}
\toprule
Metric & ASD Mean (SD) & TD Mean (SD) & Test & $p$ & $p_{\mathrm{FDR}}$ \\
\midrule
Bipolarity              & $0.591$ (0.043)   & $0.577$ (0.045)   & Mann--Whitney $U$ & $0.026$ & $0.211$ \\
Prop agreement          & $0.432$ (0.040)   & $0.442$ (0.038)   & Mann--Whitney $U$ & $0.095$ & $0.236$ \\
Prop imbalanced same    & $0.065$ (0.040)   & $0.055$ (0.038)   & Mann--Whitney $U$ & $0.109$ & $0.236$ \\
Prop imbalanced opp     & $0.433$ (0.042)   & $0.443$ (0.039)   & Mann--Whitney $U$ & $0.121$ & $0.236$ \\
Prop disagreement       & $0.070$ (0.042)   & $0.060$ (0.040)   & Mann--Whitney $U$ & $0.147$ & $0.236$ \\
Global efficiency       & $0.387$ (0.070)   & $0.402$ (0.077)   & Mann--Whitney $U$ & $0.246$ & $0.283$ \\
Local efficiency        & $0.387$ (0.070)   & $0.402$ (0.077)   & Mann--Whitney $U$ & $0.247$ & $0.283$ \\
Energy global           & $-285.4$ (146.1)  & $-308.3$ (155.7)  & Mann--Whitney $U$ & $0.305$ & $0.305$ \\
\bottomrule
\end{tabular}
}%
}
\end{table}
\clearpage

\begin{table}[H]\color{red!70!black}
\centering
\caption{Classification performance in the no-GSR sensitivity analysis (CompCor preprocessing) after
ComBat harmonization. All models performed at or near chance (cross-validated accuracy
$0.46$--$0.57$; held-out test accuracy $\leq 0.59$), consistent with the absence of FDR-significant
univariate group differences.}
\label{deletedoriginal:tab:noGSR_ml}
\small
\resizebox{\linewidth}{!}{%
\begin{tabular}{lccccc}
\toprule
Model & CV Acc & Test Acc & Test Bal Acc & Test F1 & Test AUC \\
\midrule
Gaussian Naive Bayes  & 0.567 & 0.588 & 0.594 & 0.611 & 0.531 \\
SVM (RBF)             & 0.530 & 0.500 & 0.503 & 0.514 & 0.554 \\
Logistic Regression   & 0.552 & 0.529 & 0.531 & 0.529 & 0.587 \\
KNN                   & 0.500 & 0.471 & 0.476 & 0.500 & 0.521 \\
XGBoost               & 0.455 & 0.559 & 0.559 & 0.545 & 0.556 \\
\bottomrule
\end{tabular}
}
\end{table}
\clearpage

\begin{table}[H]\color{red!70!black}
\centering
\caption{Comparison of key results across pipelines after correct site harmonization, and the
effect of harmonization itself. All group comparisons are FDR-corrected within their feature family.
``Pre-harm.'' denotes the earlier analysis in which ComBat had failed silently and raw features were
used.}
\label{deletedoriginal:tab:GSR_comparison}
\small
\resizebox{\linewidth}{!}{%
\begin{tabular}{lcc}
\toprule
\textbf{Finding} & \textbf{GSR (harmonized)} & \textbf{No-GSR / CompCor (harmonized)} \\
\midrule
Sample size (ASD / TD) & 78 / 86 & 78 / 87 \\
\midrule
\multicolumn{3}{l}{\textit{Effect of harmonization on whole-brain energy (GSR)}} \\
\quad Pre-harmonization & \multicolumn{2}{c}{$p_{\mathrm{FDR}} = 0.002$ (apparently significant)} \\
\quad Post-harmonization & \multicolumn{2}{c}{uncorrected $p = 0.048$, $p_{\mathrm{FDR}} = 0.19$ (n.s.)} \\
\midrule
\multicolumn{3}{l}{\textit{Whole-brain metrics (harmonized)}} \\
\quad Energy global & $p = 0.048$, $p_{\mathrm{FDR}} = 0.19$ (n.s.) & $p = 0.305$ (n.s.) \\
\quad Prop imbalanced same & $p = 0.047$, $p_{\mathrm{FDR}} = 0.19$ (n.s.) & $p = 0.109$ (n.s.) \\
\quad Bipolarity & $p = 0.53$ (n.s.) & $p = 0.026$, $p_{\mathrm{FDR}} = 0.21$ (n.s.) \\
\quad Graph efficiency (group diff.) & $p > 0.99$ (n.s.) & $p = 0.25$ (n.s.) \\
\midrule
\multicolumn{3}{l}{\textit{$H$--graph efficiency correlation (within ASD)}} \\
\quad Global efficiency & $r = -0.20$, $p = 0.075$ (n.s.) & $r = -0.33$, $p_{\mathrm{FDR}} = 0.005$ (sig.) \\
\midrule
\multicolumn{3}{l}{\textit{Intra-/Inter-network metrics}} \\
\quad FDR-significant effects & None & None \\
\midrule
\multicolumn{3}{l}{\textit{$H$--ADOS Total within-ASD correlation}} \\
\quad ADOS Total & $r = -0.10$, $p = 0.47$ (n.s.) & $r = +0.05$, $p = 0.71$ (n.s.) \\
\midrule
\multicolumn{3}{l}{\textit{Behavioral correlations / PLSC}} \\
\quad Any FDR-significant pair & 0 (of 460) & 0 (of 460) \\
\quad PLSC LV1 (ADOS block) & perm $p = 0.73$ (n.s.) & perm $p = 0.26$ (n.s.) \\
\midrule
\multicolumn{3}{l}{\textit{Classification (best held-out test accuracy; harmonized)}} \\
\quad Best model / accuracy & SVM / XGBoost, $0.65$ (CV $\approx 0.59$) & GaussianNB, $0.59$ (CV $\approx 0.57$) \\
\quad Interpretation & \multicolumn{2}{c}{near chance in both pipelines} \\
\bottomrule
\end{tabular}
}
\end{table}
\clearpage

\begin{table}[H]\color{red!70!black}
\centering
\caption{Demographic and sampling characteristics of the study cohort.}
\label{deletedoriginal:tab:demographics}
\setlength{\tabcolsep}{3pt}
\small
\resizebox{\columnwidth}{!}{%
\resizebox{\linewidth}{!}{%
\begin{tabular}{p{0.38\columnwidth}p{0.27\columnwidth}p{0.27\columnwidth}}
\toprule
\textbf{Characteristic} & \textbf{ASD} & \textbf{TD} \\
\midrule
Initially selected from ABIDE I, $n$ & 100 & 100 \\
Analyzed sample after QC, harmonization, and outlier removal (GSR), $n$ & 78 & 86 \\
Sex (male / female), $n$ & all male & all male \\
Age range (years) & 18--30 & 18--30 \\
Handedness & right & right \\
Full-scale IQ (FSIQ) inclusion & $>$ 80 & $>$ 80 \\
Diagnostic status & DSM ASD, ADOS/ADI-R confirmed & Typically developing (TD), no ASD diagnosis \\
Data source / preprocessing & ABIDE I, CPAC pipeline, CC200 atlas (200 ROIs) & ABIDE I, CPAC pipeline, CC200 atlas (200 ROIs) \\
Contributing sites (after excluding single/two-subject sites) & 11 & 11 \\
No-GSR (CompCor) analyzed sample (ASD / TD) & 78 / 87 & 78 / 87 \\
\bottomrule
\end{tabular}
}%
}
\begin{flushleft}
\footnotesize
All participants were right-handed adult males (18--30 years) with full-scale IQ greater than 80,
randomly selected from the ABIDE I repository and matched on age and IQ across groups. Both pipelines
were applied to the same initially selected cohort (100 ASD, 100 TD) and used identical inclusion and
motion-based quality-control criteria; the two pipelines differ only in the nuisance-regression step
(GSR versus CompCor). Prior to ComBat harmonization, two acquisition sites contributing only one
(CMU) or two (UCLA\_1) participants were excluded, as a single-participant batch cannot be
harmonized; this left 11 sites. After ComBat harmonization and Tukey IQR-based outlier removal on
whole-brain coevolutionary energy, the GSR analysis sample comprised 78 ASD and 86 TD participants,
and the no-GSR (CompCor) sample comprised 78 ASD and 87 TD participants. The small difference in
final sample size between pipelines reflects pipeline-specific outlier removal applied to each
pipeline's own energy distribution.
\end{flushleft}
\end{table}
\clearpage
\endgroup
